\documentclass[11pt]{article}
\usepackage[margin=0.3in,top=0.25in,bottom=0.4in]{geometry}
\makeatletter
\renewcommand{\@maketitle}{\newpage\null\vskip -3.5em\begin{center}{\LARGE \@title \par}\end{center}\par\vskip 0.2em}
\makeatother
\usepackage{amsmath,amssymb}
\usepackage{hyperref}

\title{D0-level fragmentation and nuclear (Glauber) targets in \texttt{flux2d}}
\author{}
\date{}

\begin{document}
\maketitle

\section{Overview}

This extends \texttt{flux2d} (see \texttt{flux2d\_overview.pdf}) with two
pieces of infrastructure ported from \texttt{inclusive-D0-UPC} and
\texttt{D0CS}:

\begin{itemize}
\item \textbf{Fragmentation}: convolving the bare charm-quark exclusive and
diffractive cross sections with a $c\to D^0$ fragmentation function, giving
$D^0$-level spectra differential in the \emph{observed meson's} transverse
momentum $p_{D^0}$ instead of the charm quark's own $p_T$. Three
fragmentation functions are available: BCFY, Kniehl--Kramer, LHAPDF.
\item \textbf{Nucleus targets}: looping the calculation over
Glauber-sampled dipole amplitude files to build the full nuclear-target
(Pb or Au) cross section, via the same impact-parameter ($b$) integration
used by \texttt{inclusive-D0-UPC}'s \texttt{run\_many\_Pb.sh} /
\texttt{cross\_section.py}.
\end{itemize}

Unlike either source repository, here \emph{both} the exclusive and the
diffractive channel get fragmented — neither D0CS nor inclusive-D0-UPC
has a fragmented exclusive (no-gluon dijet) calculation.

\section{Fragmentation}

The observed $D^0$ carries momentum fraction $z_h$ of the charm quark that
actually scatters, so the charm quark's transverse momentum is
\[
p_c = \frac{p_{D^0}}{z_h}, \qquad z_h\in[0.05,\,1.0],
\]
and $z_h$ is integrated out as an extra VEGAS dimension. The convolution
weight is $D(z_h)/z_h^2$ (\texttt{fragmentation\_weight()},
\texttt{src/integrand.cpp}), with $D(z_h)$ one of:

\begin{itemize}
\item \textbf{BCFY} (\texttt{Dc\_to\_D0()}, \texttt{src/fragmentation.cpp}):
a pseudoscalar + vector-feeddown parametrization with non-perturbative
parameter $r=0.1$.
\item \textbf{Kniehl--Kramer} (\texttt{D\_kniehl\_kramer()}):
$D(z)=N\,z(1-z)^2/[(1-z)^2+\epsilon z]^2$, $N=0.694$, $\epsilon=0.101$.
\item \textbf{LHAPDF} (\texttt{MakeLHAPDFZInterpolator()},
\texttt{src/lhapdf\_grid.cpp}): reads an LHAPDF lhagrid1-format member file
(default \texttt{data/prompt-D0-1-109/prompt-D0-1-109\_0000.dat}, member 0
only — the other 100 replicas are not shipped here, unlike
inclusive-D0-UPC's full 471~MB set) and interpolates $D(z,Q{=}m_c)$.
\end{itemize}

Selected via \texttt{flux\_grid\_d0}'s optional \texttt{frag\_type}
command-line argument (\texttt{BCFY}\,|\,\texttt{KniehlKramer}\,|
\texttt{LHAPDF}, default \texttt{BCFY} if omitted) — an explicit argument
rather than an environment variable, since it's a physics selection, not a
tuning knob (see Section~\ref{sec:target} for the same reasoning applied
to \texttt{channel}). \texttt{LHAPDF\_FILE} (still an environment
variable, since it only matters when \texttt{frag\_type=LHAPDF}) overrides
the grid file path.

\subsection{Where the BCFY $D^0$ coefficients (0.168, 0.39) come from}
\label{sec:bcfy-coeffs}

\texttt{Dc\_to\_D0()} combines the raw BCFY pseudoscalar/vector
fragmentation functions $D_P(z;r)$, $D_V(z;r)$ (Braaten, Cheung, Fleming,
Yuan, Phys.\ Rev.\ D51 (1995) 4819, hep-ph/9409316, their eqs.\ 31--32)
into the physical $c\to D^0$ fragmentation function following Cacciari \&
Nason, ``Charm Cross Sections for the Tevatron Run~II''
(hep-ph/0306212), Section~2. Physically, an observed $D^0$ can come from
three sources: primary (direct) $c\to D^0$, feed-down from $c\to D^{*0}$
followed by $D^{*0}\to D^0\,(\pi^0/\gamma)$, and feed-down from
$c\to D^{*+}$ followed by $D^{*+}\to D^0\pi^+$:
\[
F(c\to D^0) = F_p(c\to D^0) + F(c\to D^{*0})\otimes F(D^{*0}\to D^0)
                             + F(c\to D^{*+})\otimes F(D^{*+}\to D^0).
\]
Cacciari \& Nason assume isospin symmetry for the \emph{production} step
— $c\to D^{*0}$ and $c\to D^{*+}$ share the same fragmentation function
and are treated as kinematically degenerate, i.e.\ a single representative
$D^*$ mass is used for both (the actual $m_{D^{*0}}=2.0067$~GeV vs.\
$m_{D^{*+}}=2.010$~GeV splitting, $\sim$0.2\%, is neglected as below the
precision of the model) — with all of the flavor dependence pushed into
the measured \emph{decay} branching ratios. The $D^*\to D^0$ decay is
modeled as a rescaling of the momentum fraction by $m_D/m_{D^*}$,
\[
\tilde D_V(z;r) = \Theta\!\left(\frac{m_D}{m_{D^*}}-z\right)
                   D_V\!\left(\frac{m_{D^*}}{m_D}z;\,r\right)\frac{m_{D^*}}{m_D},
\]
(this is \texttt{BCFY\_DV(z\_scaled,r)*mass\_ratio} gated by
\texttt{theta\_step()} in \texttt{src/fragmentation.cpp}), and because
both decay chains share this same rescaled shape, their branching
fractions simply \emph{add} inside one overall coefficient rather than
needing two separate mass-dependent terms:
\[
D^{c\to D^0}(z;r) = \mathrm{BR}(c\to D^0_p)\,D_P(z;r)
   + \mathrm{BR}(c\to D^*)\big[\mathrm{BR}(D^{*0}\to D^0)
   + \mathrm{BR}(D^{*+}\to D^0)\big]\,\tilde D_V(z;r).
\]
With $\mathrm{BR}(c\to D^0_p)=0.168$, $\mathrm{BR}(c\to D^*)=0.233$,
$\mathrm{BR}(D^{*0}\to D^0)=1$, $\mathrm{BR}(D^{*+}\to D^0)=0.677$ (from
ALEPH, hep-ph/0306212 eqs.\ 1 and 4):
\[
0.233\times(1+0.677) = 0.390,
\]
giving exactly the $0.168$/$0.39$ pair hard-coded in
\texttt{Dc\_to\_D0()}. So the \texttt{0.39} coefficient already sums
\emph{both} the $D^{*0}\to D^0$ and $D^{*+}\to D^0$ feed-down channels —
the single \texttt{mDstar} mass in the code is not literally
``the $D^{*0}$ only'', it is the isospin-averaged $D^*$ mass standing in
for both charge states under the degenerate-mass approximation above.

\subsection{New integrands}

\texttt{integrand\_exclusive\_frag()} and
\texttt{integrand\_diffractive\_xpo\_frag()} (\texttt{src/integrand.cpp})
are self-contained siblings of \texttt{integrand\_exclusive()} /
\texttt{integrand\_diffractive\_xpo()} — not thin wrappers around them —
because $p_c$ now varies \emph{within} a single VEGAS call (once $z_h$ is
sampled), so the $H^2$-related quantities ($p_c^2+m^2$, etc.) must be
recomputed fresh on every call rather than read from the parent
integrands' precomputed, fixed-$p_T$ struct fields.

\textbf{Exclusive+frag}: 5D VEGAS over $\{r_1,r_2,u_{qp},b,z_h\}$
(\texttt{exclusiveCrossSection\_frag()}). Identical physics to the bare
exclusive kernel, evaluated at $p_c$ instead of $p_T$.

\textbf{Diffractive+frag}: 7D VEGAS over
$\{r_1,r_2,u_k,u_{qp},b,u_{x_{\mathbb P}},z_h\}$
(\texttt{diffractiveCrossSection\_xpo\_frag()}).

\subsection{Why the gluon-$k$ cutoff has to be applied differently}

The bare diffractive integrand caps the gluon momentum at $k\le p_T$ via a
hard VEGAS box edge, since $p_T$ is fixed for the whole call (see
\texttt{flux2d\_overview.pdf}, ``Diffractive process''). That specific
\emph{mechanism} can't be reused here: $p_c=p_{D^0}/z_h$ varies per
sample, so a VEGAS box edge that depends on another sampled dimension
isn't representable (VEGAS needs one fixed hyper-rectangle for the whole
call). The physical cutoff itself is unchanged though — $k\le p_c$ still
holds exactly, just enforced differently:
\begin{enumerate}
\item $k$'s VEGAS box is set to the \emph{widest} range reachable at this
$p_{D^0}$: $k\in[k_{\rm lo},\,k_{\rm upper,frag}]$ with $k_{\rm lo}=0.01$
and $k_{\rm upper,frag}=p_{D^0}/z_{\min}$ — the largest $p_c$ possible
over the whole $z_h$ range (since $p_c=p_{D^0}/z_h$ is largest at the
smallest $z_h$). This is computed once per call to
\texttt{diffractiveCrossSection\_xpo\_frag()} (like \texttt{par->log\_k}
in the bare case), not once globally, since it depends on $p_{D^0}$.
\item The exact constraint $k\le p_c$ is then applied as an explicit
per-sample check inside the integrand (\texttt{if (k > pc) return 0.0;}),
the same way other kinematic bounds (e.g.\ $0<x<1$) are enforced elsewhere
in this file.
\end{enumerate}
These are \texttt{par->k\_upper\_frag} / \texttt{par->log\_k\_frag} in
\texttt{def.hpp} — the fragmented analogs of the bare integrand's
\texttt{par->k\_upper} / \texttt{par->log\_k}, just recomputed on every
call instead of assumed fixed. (An earlier version of this note described
a Gaussian UV-suppression scheme, matching D0CS/inclusive-D0-UPC's
approach; that was replaced by this exact cutoff, which needs no arbitrary
suppression-scale parameter.)

\section{Nucleus targets (Glauber averaging)}
\label{sec:target}

The dipole amplitude $N(r,x)$ describes scattering off a specific target;
for a nucleus, different \texttt{data/<Pb,Au>/mve/glauber\_mve\_<b>} files
are independent Glauber samples at different target-nucleus impact
parameters $b$ (a different physical $b$ from the photon-flux impact
parameter already integrated inside the cross sections themselves). The
full nuclear cross section needs the additional 2D transverse-plane
integral $\int d^2b = 2\pi\int b\,db$ over these samples.

\subsection{\texttt{flux\_grid\_d0}: one point at a time}

\texttt{src/main\_grid\_d0.cpp} (built as \texttt{flux\_grid\_d0}) is
deliberately simple, matching \texttt{inclusive-D0-UPC}'s
\texttt{src/main.cpp} exactly: it computes both fragmented cross sections
at \emph{one} $(p_{D^0},y)$ point, for \emph{one} dipole file, and prints
the result — no grid, no OpenMP:
\[
\texttt{./flux\_grid\_d0 <dipole\_file> <pD0> <y> [frag\_type] [channel]}
\]
\texttt{frag\_type} and \texttt{channel} are optional positional arguments
(defaulting to \texttt{BCFY} and \texttt{An0n} if omitted) rather than
environment variables — both are physics selections that change which
formula gets used inside the integrand, so they're spelled out explicitly
on the command line rather than picked up implicitly from the shell
environment; \texttt{CALLS} and \texttt{LHAPDF\_FILE} remain environment
variables, since those are tuning knobs (how many VEGAS samples, which
grid file) rather than a choice between physical formulas. Output is a
few \texttt{\#}-prefixed header lines (fragmentation type, channel,
dipole file) followed by one data line,
\texttt{pD0 exclusive diffractive}. All of the looping over many
$(p_{D^0},y)$ points and many nucleus-target Glauber samples ($b$) —
previously done inside this program with an OpenMP grid — is now entirely
\texttt{run\_many\_nucleus.sh}'s job, described next.

\subsection{\texttt{run\_many\_nucleus.sh}: the loops and file bookkeeping}

\[
\texttt{NUCLEUS=Pb ./run\_many\_nucleus.sh}
\]
runs \texttt{flux\_grid\_d0} once for every combination of dipole file
(every \texttt{data/\$NUCLEUS/mve/glauber\_mve\_<b>} sample), $p_{D^0}$
(a plain linear sweep, \texttt{PT\_MIN} to \texttt{PT\_MAX} in steps of
\texttt{PT\_STEP} — no log-spacing trick, unlike the old internal grid),
and $y$ (\texttt{Y\_VALS}) — the same three nested \texttt{bash} loops as
\texttt{inclusive-D0-UPC}'s \texttt{run\_many\_Pb.sh}, backgrounded up to
\texttt{CORES} at once. This script is the only place \texttt{frag\_type}
and \texttt{channel} get decided (from its own \texttt{FRAG\_TYPE}/
\texttt{CHANNEL} environment variables, same defaults as
\texttt{flux\_grid\_d0}) — every invocation is then handed them explicitly
as its 4th/5th command-line arguments.

Unlike \texttt{run\_many\_Pb.sh} (and an earlier version of this script),
no temp files are used. Because \texttt{frag\_type}/\texttt{channel} (and
hence the filename tags built from them) are already known before any
point has run, every output file's header —
\texttt{files/flux\_grid\_d0\_<process>\_<frag>\_<channel>\_<NUCLEUS>\_y<Y>.dat},
rows \texttt{b pD0 dsigma} — can be written once, upfront. Each point then
captures \texttt{flux\_grid\_d0}'s output straight into a shell variable,
pulls out column 2 (exclusive) / column 3 (diffractive) with \texttt{awk},
and appends one short line directly to the corresponding file. This is
safe with many points running at once: appends under $\sim$4KB are atomic
on Linux, so concurrent points can't interleave or corrupt each other's
lines — they just can't control what \emph{order} the lines land in, so
unlike \texttt{run\_many\_Pb.sh}'s final \texttt{sort} pass, these files
come out with rows in whatever order their jobs finished, not sorted by
$b$ then $p_{D^0}$. That's immaterial downstream:
\texttt{combine\_nucleus.py} sorts by $b$ itself before integrating, so
the output file format and everything past it is unchanged.

\subsection{\texttt{combine\_nucleus.py}: the $b$ integral}

Reads all accumulated per-$y$ files for a given
(\texttt{NUCLEUS}, \texttt{FRAG\_TYPE}, \texttt{CHANNEL}), and for each
$p_{D^0}$ integrates $b\cdot d\sigma(b)$ over $b$ via Simpson's rule
(scipy), then applies:
\[
\frac{d\sigma}{dy\,dp_{D^0}} = 2\pi\Big[\!\int\! b\,d\sigma(b)\,db\Big]_{\rm Simpson}
\times \text{(prefactor)} \times 2\pi p_{D^0} \times \text{(GeV}^{-2}\to\text{mb)}.
\]
The two $2\pi$ factors are geometrically distinct: one is the Glauber-$b$
angular integral, the other is the standard $d^2p_{D^0}\to dp_{D^0}$
Jacobian (matching \texttt{plot\_flux\_grid.py}'s convention). The
per-process prefactor also matches \texttt{plot\_flux\_grid.py} exactly
(exclusive: constant; diffractive: running $\alpha_s(m_T)$) —
\emph{not} \texttt{cross\_section.py}'s \texttt{factor\_A}, which
normalizes a different physical channel (the momentum-space/S$_k$
``inclusive'' calculation that inclusive-D0-UPC computes, not ported
here). Writes \texttt{files/combined\_d0\_<process>\_...\_y<Y>.dat} and
one summary plot, \texttt{plots/combined\_d0\_<frag>\_<channel>\_<nucleus>.png}.

\section{Data files added}

\begin{itemize}
\item \texttt{data/prompt-D0-1-109/prompt-D0-1-109\_0000.dat} (+ \texttt{.info}):
central-member LHAPDF grid only (4.9~MB) — the full 101-replica set
(471~MB in \texttt{inclusive-D0-UPC}) was not copied; \texttt{LHAPDF\_FILE}
can point elsewhere if a replica band is ever needed.
\item \texttt{data/\{Pb,Au,proton\}/} were already present in
\texttt{flux2d/data} (copied from \texttt{exclusive/data} when this
repository was created) and already include the Glauber-sampled
\texttt{glauber\_mve\_<b>} files used above.
\end{itemize}

\end{document}
